\textbf{Phần 1. Trắc nghiệm lựa chọn}
\setcounterpageref{ex}{0}

\begin{ex} %[5H5V27]
Một hình tròn có chu vi là $31{,}4\text{ cm}$. Diện tích của hình tròn đó là
\choice
{$314\text{ cm}^2$}
{\True $78{,}5\text{ cm}^2$}
{$157\text{ cm}^2$}
{$62{,}8\text{ cm}^2$}
\loigiai{
Bán kính hình tròn là:
$31{,}4 : 3{,}14 : 2 = 5\text{ (cm)}$.\\
Diện tích hình tròn là:
$3{,}14 \times 5 \times 5 = 78{,}5\text{ (cm}^2\text{)}$.\\
\textbf{Đáp án của câu hỏi là B}
}
\end{ex}

\begin{ex} %[5H5V27]
Bánh xe của một xe đạp có đường kính $0{,}65\text{ m}$. Khi bánh xe lăn được đúng $100$ vòng trên mặt đất thì xe đạp đã đi được quãng đường dài bao nhiêu mét?
\choice
{$20{,}41\text{ m}$}
{\True $204{,}1\text{ m}$}
{$408{,}2\text{ m}$}
{$130\text{ m}$}
\loigiai{
Chu vi của bánh xe là:
$0{,}65 \times 3{,}14 = 2{,}041\text{ (m)}$.\\
Quãng đường xe đạp đi được sau $100$ vòng là:
$2{,}041 \times 100 = 204{,}1\text{ (m)}$.\\
\textbf{Đáp án của câu hỏi là B}
}
\end{ex}

\begin{ex} %[5H5V27]
Một bồn hoa hình tròn có bán kính $3\text{ m}$. Người ta xây một lối đi xung quanh bồn hoa rộng $1\text{ m}$. Diện tích của lối đi đó là
\choice
{$12{,}56\text{ m}^2$}
{\True $21{,}98\text{ m}^2$}
{$28{,}26\text{ m}^2$}
{$50{,}24\text{ m}^2$}
\loigiai{
Bán kính hình tròn lớn (gồm cả lối đi) là:
$3 + 1 = 4\text{ (m)}$.\\
Diện tích hình tròn lớn là:
$3{,}14 \times 4 \times 4 = 50{,}24\text{ (m}^2\text{)}$.\\
Diện tích bồn hoa ban đầu là:
$3{,}14 \times 3 \times 3 = 28{,}26\text{ (m}^2\text{)}$.\\
Diện tích của lối đi là:
$50{,}24 - 28{,}26 = 21{,}98\text{ (m}^2\text{)}$.\\
\textbf{Đáp án của câu hỏi là B}
}
\end{ex}

\begin{ex} %[5H5V27]
Nếu tăng bán kính của một hình tròn lên gấp 2 lần thì chu vi và diện tích của hình tròn đó thay đổi như thế nào?
\choice
{Chu vi gấp 4 lần, diện tích gấp 4 lần}
{Chu vi gấp 2 lần, diện tích gấp 2 lần}
{\True Chu vi gấp 2 lần, diện tích gấp 4 lần}
{Chu vi gấp 4 lần, diện tích gấp 2 lần}
\loigiai{
Chu vi hình tròn: $C = 2 \times 3{,}14 \times r$, nên khi $r$ gấp 2 lần thì chu vi gấp 2 lần.\\
Diện tích hình tròn: $S = 3{,}14 \times r \times r$, nên khi $r$ gấp 2 lần thì diện tích gấp $2 \times 2 = 4$ lần.\\
\textbf{Đáp án của câu hỏi là C}
}
\end{ex}

\begin{ex} %[5H5V27]
Một mặt bàn hình tròn có chu vi đo được $3{,}768\text{ m}$. Diện tích của mặt bàn đó là
\choice
{$11{,}304\text{ m}^2$}
{\True $1{,}1304\text{ m}^2$}
{$4{,}5216\text{ m}^2$}
{$2{,}2608\text{ m}^2$}
\loigiai{
Bán kính mặt bàn là:
$3{,}768 : 3{,}14 : 2 = 0{,}6\text{ (m)}$.\\
Diện tích mặt bàn là:
$3{,}14 \times 0{,}6 \times 0{,}6 = 1{,}1304\text{ (m}^2\text{)}$.\\
\textbf{Đáp án của câu hỏi là B}
}
\end{ex}

\begin{ex} %[5H5V27]
Một người đi xe đạp đi được quãng đường dài $471\text{ m}$ khi bánh xe lăn được $250$ vòng. Đường kính của bánh xe đạp đó là
\choice
{$50\text{ cm}$}
{\True $60\text{ cm}$}
{$70\text{ cm}$}
{$80\text{ cm}$}
\loigiai{
Chu vi của bánh xe là:
$471 : 250 = 1{,}884\text{ (m)}$.\\
Đường kính của bánh xe là:
$1{,}884 : 3{,}14 = 0{,}6\text{ (m)} = 60\text{ cm}$.\\
\textbf{Đáp án của câu hỏi là B}
}
\end{ex}

\begin{ex} %[5H5V27]
Một tấm bìa hình vuông có cạnh $8\text{ cm}$. Người ta vẽ một hình tròn nội tiếp tiếp xúc với 4 cạnh của hình vuông đó (đường kính hình tròn bằng $8\text{ cm}$) rồi cắt bỏ phần hình tròn đi. Diện tích phần bìa còn lại là
\choice
{$12{,}56\text{ cm}^2$}
{\True $13{,}76\text{ cm}^2$}
{$14{,}24\text{ cm}^2$}
{$16\text{ cm}^2$}
\loigiai{
Diện tích tấm bìa hình vuông là:
$8 \times 8 = 64\text{ (cm}^2\text{)}$.\\
Bán kính hình tròn là:
$8 : 2 = 4\text{ (cm)}$.\\
Diện tích hình tròn là:
$3{,}14 \times 4 \times 4 = 50{,}24\text{ (cm}^2\text{)}$.\\
Diện tích phần bìa còn lại là:
$64 - 50{,}24 = 13{,}76\text{ (cm}^2\text{)}$.\\
\textbf{Đáp án của câu hỏi là B}
}
\end{ex}

\begin{ex} %[5H5V27]
Hai hình tròn có hiệu hai bán kính bằng $2\text{ cm}$. Hiệu hai chu vi của hai hình tròn đó là
\choice
{$6{,}28\text{ cm}$}
{\True $12{,}56\text{ cm}$}
{$18{,}84\text{ cm}$}
{$25{,}12\text{ cm}$}
\loigiai{
Hiệu chu vi của hai hình tròn là:
$C_1 - C_2 = 2 \times 3{,}14 \times r_1 - 2 \times 3{,}14 \times r_2 = 2 \times 3{,}14 \times (r_1 - r_2) = 6{,}28 \times 2 = 12{,}56\text{ (cm)}$.\\
\textbf{Đáp án của câu hỏi là B}
}
\end{ex}

\begin{ex} %[5H5V27]
Một mảnh đất có dạng hình bán nguyệt (nửa hình tròn) với đường kính đáy $20\text{ m}$. Chu vi của mảnh đất đó là
\choice
{$31{,}4\text{ m}$}
{\True $51{,}4\text{ m}$}
{$62{,}8\text{ m}$}
{$41{,}4\text{ m}$}
\loigiai{
Độ dài cung tròn nửa đường tròn là:
$20 \times 3{,}14 : 2 = 31{,}4\text{ (m)}$.\\
Chu vi mảnh đất hình bán nguyệt gồm độ dài cung tròn cộng với độ dài đường kính:
$31{,}4 + 20 = 51{,}4\text{ (m)}$.\\
\textbf{Đáp án của câu hỏi là B}
}
\end{ex}

\begin{ex} %[5H5V27]
Nếu diện tích của một hình tròn tăng lên gấp 9 lần thì bán kính của hình tròn đó
\choice
{Tăng gấp 9 lần}
{\True Tăng gấp 3 lần}
{Tăng gấp 4{,}5 lần}
{Tăng gấp 6 lần}
\loigiai{
Diện tích hình tròn $S = 3{,}14 \times r \times r$.\\
Để diện tích tăng lên $9$ lần thì tích $r \times r$ phải tăng lên $9$ lần.\\
Vì $3 \times 3 = 9$ nên bán kính $r$ phải tăng lên gấp $3$ lần.\\
\textbf{Đáp án của câu hỏi là B}
}
\end{ex}

\begin{ex} %[5H5V27]
Bác Nam dùng một sợi dây thép quấn quanh một cột trụ bê tông hình tròn có đường kính $0{,}8\text{ m}$ vừa vặn đúng 5 vòng. Chiều dài sợi dây thép bác Nam đã dùng là
\choice
{$6{,}28\text{ m}$}
{\True $12{,}56\text{ m}$}
{$25{,}12\text{ m}$}
{$10\text{ m}$}
\loigiai{
Chu vi một vòng quanh cột trụ là:
$0{,}8 \times 3{,}14 = 2{,}512\text{ (m)}$.\\
Chiều dài sợi dây thép đã quấn 5 vòng là:
$2{,}512 \times 5 = 12{,}56\text{ (m)}$.\\
\textbf{Đáp án của câu hỏi là B}
}
\end{ex}

\begin{ex} %[5H5V27]
Cho hai hình tròn đồng tâm có bán kính lần lượt là $5\text{ cm}$ và $7\text{ cm}$. Diện tích của phần hình vành khăn nằm giữa hai đường tròn đó là
\choice
{$37{,}68\text{ cm}^2$}
{\True $75{,}36\text{ cm}^2$}
{$153{,}86\text{ cm}^2$}
{$78{,}5\text{ cm}^2$}
\loigiai{
Diện tích hình tròn lớn là:
$3{,}14 \times 7 \times 7 = 153{,}86\text{ (cm}^2\text{)}$.\\
Diện tích hình tròn nhỏ là:
$3{,}14 \times 5 \times 5 = 78{,}5\text{ (cm}^2\text{)}$.\\
Diện tích hình vành khăn là:
$153{,}86 - 78{,}5 = 75{,}36\text{ (cm}^2\text{)}$.\\
\textbf{Đáp án của câu hỏi là B}
}
\end{ex}

\textbf{Phần 2. Trắc nghiệm đúng sai}
\setcounterpageref{ex}{0}

\begin{ex} %[5H5V27]
Cho một hình tròn có bán kính $r = 4\text{ cm}$.
\choiceTF
{\True Đường kính của hình tròn là $8\text{ cm}$}
{\True Chu vi của hình tròn là $25{,}12\text{ cm}$}
{\True Diện tích của hình tròn là $50{,}24\text{ cm}^2$}
{Khi tăng bán kính lên $3$ lần thì diện tích của hình tròn tăng lên $6$ lần}
\loigiai{
\begin{itemchoice}
\item Đúng. Đường kính $d = 2 \times 4 = 8\text{ (cm)}$.
\item Đúng. Chu vi $C = 8 \times 3{,}14 = 25{,}12\text{ (cm)}$.
\item Đúng. Diện tích $S = 3{,}14 \times 4 \times 4 = 50{,}24\text{ (cm}^2\text{)}$.
\item Sai. Khi bán kính tăng 3 lần thì diện tích tăng $3 \times 3 = 9$ lần.
\end{itemchoice}
}
\end{ex}

\begin{ex} %[5H5V27]
Một bánh xe máy có đường kính $0{,}5\text{ m}$.
\choiceTF
{\True Chu vi của bánh xe máy là $1{,}57\text{ m}$}
{\True Khi bánh xe quay được $100$ vòng thì xe máy đi được quãng đường $157\text{ m}$}
{\True Để đi được quãng đường $1\,570\text{ m}$ thì bánh xe phải quay đúng $1\,000$ vòng}
{Nếu thay bằng bánh xe có đường kính lớn hơn thì với cùng $100$ vòng quay, quãng đường đi được sẽ ngắn hơn}
\loigiai{
\begin{itemchoice}
\item Đúng. Chu vi: $0{,}5 \times 3{,}14 = 1{,}57\text{ (m)}$.
\item Đúng. Quãng đường: $1{,}57 \times 100 = 157\text{ (m)}$.
\item Đúng. Số vòng quay: $1\,570 : 1{,}57 = 1\,000\text{ (vòng)}$.
\item Sai. Bánh xe có đường kính lớn hơn thì chu vi lớn hơn, đi được quãng đường dài hơn.
\end{itemchoice}
}
\end{ex}

\textbf{Phần 3. Trắc nghiệm trả lời ngắn}
\setcounterpageref{ex}{0}

\begin{ex} %[5H5V27]
Một bánh xe đạp có chu vi $2\text{ m}$. Hỏi bánh xe phải quay bao nhiêu vòng để xe đi được quãng đường dài $500\text{ m}$?
\shortans{250}
\loigiai{
Số vòng quay cần thiết là:
$500 : 2 = 250\text{ (vòng)}$.
}
\end{ex}

\begin{ex} %[5H5V27]
Một hình tròn có diện tích bằng $314\text{ cm}^2$. Tính đường kính của hình tròn đó theo đơn vị xăng-ti-mét.
\shortans{20}
\loigiai{
Tích hai bán kính là:
$r \times r = 314 : 3{,}14 = 100$.\\
Vì $10 \times 10 = 100$ nên bán kính hình tròn là $10\text{ cm}$.\\
Đường kính của hình tròn là:
$10 \times 2 = 20\text{ (cm)}$.
}
\end{ex}

\begin{ex} %[5H5V27]
Bán kính hình tròn lớn gấp $3$ lần bán kính hình tròn bé. Hỏi diện tích hình tròn lớn gấp bao nhiêu lần diện tích hình tròn bé?
\shortans{9}
\loigiai{
Diện tích hình tròn tỉ lệ thuận với bình phương bán kính.\\
Khi bán kính gấp 3 lần thì diện tích gấp:
$3 \times 3 = 9\text{ (lần)}$.
}
\end{ex}

\begin{ex} %[5H5V27]
Một biển báo giao thông hình tròn có bán kính $30\text{ cm}$ (diện tích là $2\,826\text{ cm}^2$). Ở giữa có một vạch ngang hình chữ nhật màu trắng kích thước $10\text{ cm} \times 40\text{ cm}$, phần còn lại sơn màu đỏ. Tính diện tích phần sơn màu đỏ theo đơn vị xăng-ti-mét vuông.
\shortans{2426}
\loigiai{
Diện tích vạch chữ nhật màu trắng là:
$10 \times 40 = 400\text{ (cm}^2\text{)}$.\\
Diện tích phần sơn màu đỏ là:
$2\,826 - 400 = 2\,426\text{ (cm}^2\text{)}$.
}
\end{ex}

\begin{ex} %[5H5V27]
Một sợi dây thép uốn thành một hình tròn có bán kính $50\text{ cm}$. Tính chiều dài sợi dây thép đó theo đơn vị xăng-ti-mét.
\shortans{314}
\loigiai{
Chiều dài sợi dây thép chính bằng chu vi hình tròn:
$2 \times 3{,}14 \times 50 = 314\text{ (cm)}$.
}
\end{ex}

\begin{ex} %[5H5V27]
Một hồ nước hình tròn có diện tích đo được là $1\,256\text{ m}^2$. Tính đường kính của hồ nước đó theo đơn vị mét.
\shortans{40}
\loigiai{
Tích hai bán kính của hồ nước là:
$r \times r = 1\,256 : 3{,}14 = 400$.\\
Vì $20 \times 20 = 400$ nên bán kính hồ nước là $20\text{ m}$.\\
Đường kính của hồ nước là:
$20 \times 2 = 40\text{ (m)}$.
}
\end{ex}
